% Volume-owned synthesis, status, roadmap, concordance, and bibliography.
% Kept separate from the orchestration root so the publication remains a
% bounded two-chunk review even after the compiled PDF is attached.

\clearpage
\phantomsection\label{book:appendices}
\addcontentsline{toc}{part}{\textbf{\scshape Appendices}}
\begin{center}
  \vspace*{0.3cm}
  {\Huge\scshape\bfseries Appendices\par}
  \vspace{0.2cm}
  {\color{hhcobalt}\rule{0.3\textwidth}{1.5pt}}\par
  \vspace{0.4cm}
\end{center}
\appendix
\renewcommand{\thesection}{\Alph{section}}
\renewcommand{\thefigure}{A.\arabic{figure}}
\renewcommand{\thetable}{A.\arabic{table}}
\renewcommand{\theHsection}{volume.appendix.\arabic{section}}
\renewcommand{\theHsubsection}{volume.appendix.\arabic{section}.\arabic{subsection}}
\renewcommand{\theHsubsubsection}{volume.appendix.\arabic{section}.\arabic{subsection}.\arabic{subsubsection}}
\renewcommand{\theHfigure}{volume.appendix.\arabic{figure}}
\renewcommand{\theHtable}{volume.appendix.\arabic{table}}
\renewcommand{\theHequation}{volume.appendix.\arabic{equation}}
\renewcommand{\theHlstlisting}{volume.appendix.\arabic{lstlisting}}
\pdbackmatterheaders{Appendices}

\section{Implementation boundary}\label{app:volume-status}

The book argues for more than the product does today. This appendix says, for
each mechanism, which parts run now and which exist only as a design or a
proof, so that neither is read as the other. A mechanism gets the grade of its
weakest necessary part. A verified library does not make every program that
calls it safe, and offering a checked route does not show that every effect
takes it.

\subsection{Five assurance modes}
\pdchapref{swk}{The Single-Writer Kernel} defines the five assurance modes
where it first uses them (\S\ref{swk:sec:partial-observation}). This list
repeats them for reference.
% Recut from leftmargin=2.6cm,style=nextline (set for the book's earlier A4
% trim -- more than half the 4.5in measure) to a run-in head: bold term, same
% line as its text, small hanging indent in em.
\begin{description}[leftmargin=1.4em,itemsep=3pt]
  \item[Observed] The act may happen outside the system; the system keeps
  best-effort evidence after the fact.
  \item[Coordinated] Cooperative clients register claims, roles, commitments,
  and messages before acting.
  \item[Brokered] A chosen credential or effect channel passes through a broker
  that the daemon owns.
  \item[Confined] Operating-system controls leave the process no path to the
  protected effect other than the paths it was given.
  \item[Attested] A measured environment releases a key, or accepts a request,
  only for a runtime whose identity it can verify.
\end{description}

The modes are not stages of maturity. They are not evenly spaced, and they do
not even form one order, because Brokered and Confined stop different
attackers. Each mode names what an attacker would have to defeat to make a
claim about the mechanism false. At Observed that is nothing; at Coordinated,
a client's willingness to register first; at Brokered, a broker the daemon
owns; at Confined, the operating system; at Attested, a measured environment
and a key released only to it.

\begin{pdexample}{One Write, Five Modes}
\setlength{\parskip}{4pt plus 1pt}
This is a constructed example. Heron, the agent repairing the checkout test,
wants to write to the staging database without going through the daemon. It
has found the database password in an environment variable. Run that one
attempt through each mode and ask whether the claim ``every write to staging
was checked'' survives.

\emph{Observed.} The write succeeds. The operator sees it later, in the
database's own log.

\emph{Coordinated.} The write succeeds. Heron never registers it, and nothing
in a protocol for cooperative clients makes it register.

\emph{Brokered.} The write succeeds. The broker holds the daemon's credential,
but the leaked password is a second credential that never passes through the
broker.

\emph{Confined.} The write fails. The operating system gives Heron's process no
network path to the database except the one through the broker, so the
password has nowhere to go.

\emph{Attested.} The write fails. The database accepts writes only from an
environment whose measurement it has verified, and Heron's process does not
run in one.

The claim survives in the last two modes, for different reasons. Revoking the
leaked password would have made the Brokered claim hold. It would change
nothing at Confined, which never depended on Heron lacking a password.
\end{pdexample}

The modes describe one claim at run time by what it rests on. Automation
levels describe how work is divided between a person and a machine, and
evaluation assurance levels describe how closely a security claim was
examined. Chapter~1 compares the three. A system can be more automated without
a stronger guard on its effects, and more thoroughly evaluated without
dividing its decisions differently.

\subsection{What exists now}

\paragraph{Local transactions.}
\pdstatus{\BuiltWeak}{}
The daemon writes sessions, claims, messages, commitments and some other
coordination state through one SQLite writer in write-ahead-log mode, and
regression tests exercise those paths. Durability guarantees still depend on
the kind of failure; a process crash and a power loss are covered differently.
Schema
changes do not yet follow the strongest design in \pdchapref{swk}{The
Single-Writer Kernel}, and file or network effects that bypass the daemon's
checked paths are outside the transaction.

\paragraph{Operator views.}
\pdstatus{\BuiltWeak}{}
The product builds the operator's summary views from the event and
coordination records, and each item links back to the records it came from.
No view can beat the minimum-information limit proved in \pdchapref{ls}{The
Legible Swarm}, so each view's coverage, what it leaves out, and how long a
reviewer takes to see an event still have to be measured.

\paragraph{Durable work records.}
\pdstatus{\BuiltWeak}{}
Commitments, deadlines, messages, outcomes and recovery evidence survive the
end of a process. Two parts are designed, not built: a package that lets a job
resume unchanged on any machine, and a way to restore a job's working state
that does not depend on one model provider.

\paragraph{Authentication and attenuation.}
\pdstatus{\BuiltWeak}{protocol model proved; code checked on bounded inputs}
ProVerif, an automatic protocol verifier, proves the recorded
proof-of-possession and every-hop attenuation queries for the protocol model.
The Kani checks of the deployed Rust crate are bounded: the parser does not
panic on 32-byte inputs with the cryptography stubbed out, the comparator does
not panic, and two concrete vectors pass the subset check. The code that calls
the primitive must still be shown to conform, which means showing that every
protected effect is actually checked by the verified primitive.

\paragraph{Sealed execution.}
\pdstatus{\Designed}{properties model-checked}
The sealed room of \pdchapref{sealed}{The Sealed Harbor} is a design whose
properties have been model-checked. Noninterference modulo declassification
and the invariants of the release ledger are checked exhaustively on finite
models; $\varepsilon$-conservation, canary power and Wald latency are theorems
over the stated models. No harbor, meaning a daemon and its database governing
a group of agents, yet runs a work order inside the two fences (a monitor
inside the guest and a gateway outside it) with one metered release slot. The
attestation, the gate and the ledger are designed, not built, so the
\emph{Confined} mode they would supply is not yet available to any caller.

\paragraph{Escrow and settlement.}
\pdstatus{\BuiltWeak}{custody and settlement designed only}
The ledger and its conservation checks are partly built. Custody across
authorities, adjudication and 2-of-3 settlement are designed, not built, and
each depends on a custody interface that no party can bypass.

\paragraph{Relay federation.}
\pdstatus{\BuiltWeak}{}
The current Relay carries authenticated events between harbors over outbound
HTTPS and server-sent events.
\pdscope{The Relay does not copy daemon databases, put events in one global
order, authorize local effects, keep messages secret end to end on its own, or
provide the witness logs and settlement rules of \pdchapref{fh}{The Federated
Harbor}.}
% Research-agenda addition (control strand): what isolation supplies.
% Inserted in Appendix A after "What exists now", before the mechanized claims.
% The appendices are not passed through the Book generator's citation
% rewriting, so sources are named in prose; their entries are cited from
% The Single-Writer Kernel (section "The control game") and appear in the
% collected references.
\subsection{What isolation supplies, and what it leaves}\label{app:isolation}

The book's guarantees hold on mediated events. In
\pdchapref{swk}{The Single-Writer Kernel},
Theorem~\ref{swk:thm:regimentation-controllability} takes $\Sigma_c$ to be the
set of events \emph{actually} mediated, and
Theorem~\ref{swk:thm:swk-control-game} shows that one unmediated exit makes
every protocol unsafe against a worst-case agent. Which events are mediated is
decided below the gate, by the isolation layer. That layer sorts an agent's
possible effects into three classes. An effect is \textbf{absent} when no
device, path or credential for it exists. It is \textbf{gated} when the path
exists and passes a gate the harbor owns. It is \textbf{residual} when the
path exists and bypasses the gate. Guarantees attach to the gated class. A
residual channel needs a capacity bound or an explicit exclusion; an absent
one costs nothing to govern.

Table~\ref{tab:app-isolation-primitives} grades six primitives against that
sorting. The sources are Agache and colleagues on Firecracker, Young and
colleagues on gVisor, Klein and colleagues and Murray and colleagues on seL4,
Lehmann, Kinder and Pradel on WebAssembly, Chuang and colleagues on TEE.fail,
Marchand and colleagues on container escape, and Debenedetti and colleagues
on CaMeL; their entries are cited from the control-game section of
\pdchapref{swk}{The Single-Writer Kernel}.

\begin{table}[H]
\centering
\caption{What each isolation primitive makes absent or gated, and what it
leaves residual. Figures are the cited authors' reported values, not
measurements of this book's reference system.}
\label{tab:app-isolation-primitives}
\small
\begin{tabularx}{\textwidth}{@{}>{\raggedright\arraybackslash}p{0.15\textwidth} >{\raggedright\arraybackslash}X >{\raggedright\arraybackslash}X@{}}
\toprule
\textbf{Primitive} & \textbf{Absent or gated} & \textbf{Residual} \\
\midrule
Containers (namespaces, seccomp) & Paths outside the namespaces; system calls the filter refuses. & The shared host kernel, runtime bugs and misconfiguration. SandboxEscapeBench finds that models identify and exploit escape routes whenever misconfigurations or privilege mistakes are present. \\
gVisor & The guest's system calls, answered by a user-space kernel: $211$ of $319$ Linux calls implemented, $55$ host calls used, \texttt{open} and \texttt{socket} excluded. & Flaws in that kernel, and timing. Cost: the fastest system-call path is $2.8\times$ slower than a native process, and opening and closing a file on an external tmpfs is $216\times$ slower than from a plain container. \\
Firecracker microVM & Every device outside a four-device model (virtio network and block, serial, part of a keyboard controller); a jailer admitting $24$ system calls; about $50{,}000$ lines of Rust against over $1.4$ million for QEMU. & Whatever crosses the virtio network device; the host kernel; microarchitectural channels unless the host disables SMT and enables mitigations; power and thermal channels. \\
seL4 & Storage channels of the abstract specification, by a machine-checked intransitive noninterference proof carried to the C code by refinement. & Covert timing channels, of which the proof ``says little''. It assumes disabled DMA, a static partition schedule, and a correct compiler and hardware. \\
WebAssembly & Host memory and code outside the module; typed, structured control flow. & Memory safety inside linear memory: no stack canaries, no guard pages, writable constants. Demonstrated attacks include cross-site scripting in a browser, remote code execution in Node.js and an arbitrary file write in wasmtime. \\
Trusted execution (TDX, SEV-SNP, GPU) & Plaintext from host software; a measurement that can condition key release. & Physical access. TEE.fail extracts attestation keys with a DDR5 bus interposer costing under $\$1{,}000$; a forged TDX quote passed Intel's verification at its ``UpToDate'' level; GPU attestation not bound to the confidential VM can be relayed. Intel and AMD treat interposition as out of scope. \\
\bottomrule
\end{tabularx}
\end{table}

Three consequences follow for the book's grades.

\paragraph{Confinement for hostile code needs a VM or microVM.}
The Confined mode above says the process cannot reach the protected effect
through an ambient path. A namespace sandbox shares the host kernel with
everything it protects, and the SandboxEscapeBench authors recommend treating
plain Docker isolation as insufficient by default. The hardening specified for
OP-9 in \pdchapref{swk}{The Single-Writer Kernel}
(\S\ref{swk:sec:merkle}), \texttt{bwrap} or \texttt{unshare} sandboxes under
ephemeral user identities, moves the adversary to the different-user row of
the threat table. For an agent that runs arbitrary code and may exploit kernel
or runtime flaws, a Confined claim needs a hypervisor boundary, with
namespaces and system-call filters as defense in depth inside it. Even a
microVM keeps its network device: confinement removes ambient paths; it does
not decide what crosses the gated ones.

\paragraph{Exfiltration over allowed tools is semantic.}
Once isolation has fixed the alphabet, an allowed channel carries whatever
data flow the policy admits. CaMeL separates a task's control flow from its
data flow, tags values with capabilities, and checks a policy at every tool
call. It solves $77\%$ of the AgentDojo tasks with provable security, against
$84\%$ for an undefended agent. It lists text-to-text attacks as a non-goal and
acknowledges side channels. That is the structure of the attenuation of
\pdchapref{anchor}{The Anchor Protocol} and of the labels and release gate of
\pdchapref{sealed}{The Sealed Harbor}. No hypervisor supplies it.

\paragraph{The model provider is a named principal.}
\pdchapref{sealed}{The Sealed Harbor} notes that a remote model interface is a
trusted reader of plaintext, and places the model inside the confidential
domain (\S\ref{sealed:sec:sealed-design}). When the parties cannot hold the
weights, the provider is a third principal and the work order must name it.
Confidential-computing inference moves that trust to the vendor's attestation
keys and to whoever has physical custody of the host. After TEE.fail, an
Attested claim must state physical custody as an assumption, because the
vendors exclude interposition from their threat models.

\paragraph{An enumeration protocol (\Vision{}).}
Run the reference worker under a namespace sandbox, under gVisor and under
Firecracker. For each, attack the channel list of the kernel chapter's
boundary item~(ii): raw sockets, DNS, inherited file descriptors,
\texttt{/proc}, credential helpers, package managers and provider callbacks.
Add escape tasks of the SandboxEscapeBench kind. Record from outside the guest
whether each channel is absent, gated or residual, and measure the capacity of
timing channels over the permitted output in bits per second. Logs written
inside the guest may corroborate; they cannot decide. The result is the
$\Sigma_c$/$\Sigma_u$ inventory that the kernel chapter says a red team, not a
proof, must supply. Nothing here has been run.


% The book's one Overfull \vbox (2.96pt) sits inside this generated table
% set, one page break into its chain of eight xltabular blocks -- not from
% too little stretch (\raggedbottom alone, tried first, changed nothing:
% ragged/flush bottom only decides how a SHORT page is padded, and this
% page's problem is a row combination that is too TALL, which neither
% setting shrinks) but from real, reclaimable, redundant space: each table
% already opens on its own \medskip from the generator, and longtable/
% xltabular's machinery then adds its OWN \LTpre skip on top of that before
% the caption -- 6pt plus 12(+-4)pt for eight tables neither of us needed
% doubled. Dropping longtable's own pre/post skip to zero here removes that
% redundancy across the whole appendix, which is what actually gives the one
% tight break room to land without overflowing. \raggedbottom stays too:
% a dense reference appendix does not need flush-bottom's page-filling
% stretch, and it costs nothing to leave a short page short instead of
% stretched.
{\raggedbottom
\setlength{\LTpre}{0pt}\setlength{\LTpost}{0pt}
\pdgeneratedinput{../.cache/whitepaper-build/coordination-papers-mega-volume/mega-volume-mechanized.tex}{mechanized claims}
\par}

\section{Result atlas}\label{app:result-atlas}\label{app:one-spine}

The research program contains derivations, experiments, literature comparisons
and counterexamples. Its claims are tracked individually. Each
result below has a stable identifier, assumptions, a boundary, a chapter home,
and a named verification artifact in the companion manifest.

\subsection{The eight propositions and their evidence}

\paragraph{1. Prevention begins at the effect boundary (R5, R8) --- \pdchapref{swk}{The Single-Writer Kernel}.}
Controllability and sufficient observation together distinguish enforceable
pre-effect rules from properties requiring other controls or post-effect remedies. The airlock experiments show why complete
channel mediation is the decisive assumption. A token, label, or policy document
does not control an effect channel by itself.

\paragraph{2. Delegation only narrows (ANCHOR-ATTENUATION, R5) --- \pdchapref{anchor}{The Anchor Protocol}.}
Every hop must authenticate the issuer and reject a capability that expands
scope or lifetime. The cryptographic models close the stated monotonicity
property; the engineering obligation is universal use at protected channels.

\paragraph{3. Confinement is a property of the room, not of the tenant (R9, R10, R11) --- \pdchapref{sealed}{The Sealed Harbor}.}
Two owners who distrust each other compute one answer inside two fences with a
single metered slot. Noninterference modulo declassification holds on the
finite room model; the release ledger conserves its $\varepsilon$ budget and
refuses the torn write; canaries detect exfiltration with a stated power curve
and a Wald clock. For $q$ fixed-length $b$-bit releases, $q\cdot b$ bounds payload capacity before
timing, and timing, cache, and physical channels are declared out of model
rather than closed.

\paragraph{4. Supervision has an information cost (R1--R4, R14, R16) --- \pdchapref{ls}{The Legible Swarm}.}
The information floor prices guaranteed recall; the dual-reader result shows why
one ordering cannot generally serve both successor selection and operator
attention; the regret rule derives intervention from asymmetric error cost;
adaptive zoom and context paging describe when attention can be spent later
rather than summarized away now. The bounds price guarantees under their stated
event and flag models. They do not certify an arbitrary dashboard.

\paragraph{5. Accountability outlives the process (R7, R12, R13, B6) --- \pdchapref{stp}{From Spawn to Person}.}
The inspection tower prices reputation as amortized verification. No-mint transfer semantics prevent skill or identity replacement from creating
new authority. Budget-pooling and exposure arguments make reset laundering a
property of the principal's case file, not the process name. The admission and exposure assumptions are essential: an authentic key alone
does not prevent its holder from abandoning liabilities under a new key.

\paragraph{6. Exchange separates payment, collateral, evidence, and closure
(R15, R17) --- \pdchapref{he}{The Harbor Economy}.}
Conservation supplies the accounting invariant; the tractable deontic fragment
decides which rule sets can conflict and names the frontier past which that
decision turns NP-complete; the succession rule has a price, because under the stated queue model a sole
owner has a residual-downtime floor; comparison with a stable pooled queue
gives a conditional pooling threshold. These are conditional economic
statements. Detection probability, one-shot gain, collateral, judge
independence, and custody must be measured or contracted rather than implied.

\paragraph{7. Uncertainty needs an institution (the $\delta^{*}$ claim-signaling
game, the conservation model; R11 as instrument) --- \pdchapref{bonded}{The Bonded Commons}.}
Atomic ledgers preserve conservation; truthful claim signaling is
incentive-compatible above the mechanized discount threshold $\delta^{*}$;
audit sampling and canaries borrow their power curve from the sealed room's
detector; any claimed benefit of randomized adjudication depends on its
selection, independence and enforcement assumptions. Evidence may support a remedy without
becoming mechanical truth.

\paragraph{8. Federation relays evidence, not sovereignty (R6, R6-CR,
RELAY-V0) --- \pdchapref{fh}{The Federated Harbor}.}
Signed roots and redundant comparison paths can expose equivocation across a
specified topology. The counterexample R6-CR rejects the tempting claim that a
zero local residual proves global consistency. Current Relay establishes event
transport and authentication; it does not close witness, secrecy, liveness,
custody, or local authorization obligations.

\subsection{A result is admitted only with its limit}
For this volume, a claim is incomplete unless it states all five of the
following: the object it concerns, the assumptions under which it holds, the
kind of evidence offered, a boundary or falsifier, and the artifact by which a
reader can recompute or inspect it. The machine-readable manifest is the
publication ledger for that discipline. New mechanisms belong in the book only
after they can name the result they instantiate and the obligation they leave
unresolved.

\section{Notation and cross-chapter concordance}\label{app:concordance}
% Recut from leftmargin=2.4cm,style=nextline for the same reason as the
% assurance-modes list above.
\begin{description}[leftmargin=1.4em,itemsep=4pt]
  \item[Work unit] The durable case file joining authorization, capability
  history, evidence, obligations, outcomes, and settlement.
  \item[Harbor] One local administrative and transactional authority, normally
  one daemon and its durable database.
  \item[Principal] The human or organization whose policy and resources an
  agent ultimately exercises; not interchangeable with a process identity.
  \item[Role] A bundle of obligations, capabilities, and authority that a spawn
  may fill.
  \item[Person] An accountable identity with continuity and a witnessed outcome history;
  roles may change without creating a new accountable principal. Its anti-reset
  properties depend on admission and liability rules.
  \item[Oracle] An external or independently checkable reference used to close
  a commitment or adjudicate an outcome. Its trust model must be named.
  \item[Epoch] A version or freshness counter, with chapter-specific meaning.
  Chapter 2 uses revocation epochs; Chapter 3 uses attestation freshness to
  detect rollback. A credential's time-to-live is a validity duration, not
  another epoch counter. None substitutes for the others.
  \item[Round] A model step in a dissemination or protocol analysis, not
  automatically a wall-clock deadline.
  \item[Closed] The exact recorded property has been mechanically discharged.
  No whole-system implication follows without a conformance argument.
\end{description}

\subsection*{Which chapter owns which primitive}
\label{app:ownership}

Each shared primitive has a canonical definition in its owning chapter.
Later chapters recall the relevant part in context and apply it to a new
example; readers should not have to retrieve a distant definition merely to
follow a paragraph. Those reminders must preserve the same contract. A genuine
extension states the added assumptions explicitly. The table identifies where
to resolve a conflict, not a prohibition on useful repetition.

\begin{table}[H]
\centering
\caption{Primitive ownership. ``Owner'' defines it and proves what is proved
about it; ``cites'' names the chapters that use it without redefining it.}
\label{tab:primitive-ownership}
\small
\begin{tabularx}{\textwidth}{@{}>{\raggedright\arraybackslash}X >{\raggedright\arraybackslash}p{0.26\textwidth} >{\raggedright\arraybackslash}p{0.15\textwidth}@{}}
\toprule
\textbf{Primitive} & \textbf{Owner} & \textbf{Cites it} \\
\midrule
Authorization chain (attenuation, hop-binding, replay) & ch.~2, Anchor Protocol & ch.~1, 5, 7, 8 \\
Coordination lineage (task graph, loop detection) & ch.~1, Single-Writer Kernel & ch.~4 \\
Leakage budget and the priced exit & ch.~3, Sealed Harbor & ch.~4, 6 \\
Float plan & ch.~6, Harbor Economy & ch.~7, 8 \\
Bond, escrow, and the commons pool & ch.~7, Bonded Commons & ch.~6, 8 \\
Cross-harbor transfer ceremony & ch.~8, Federated Harbor & ch.~6 \\
Witness log, epoch root, equivocation residual & ch.~8, Federated Harbor & ch.~2, 7 \\
Threat bands and the visibility tiers & ch.~8, Federated Harbor & ch.~3 \\
Keystone identity (non-forgeable principal) & ch.~5, From Spawn to Person & ch.~6, 7 \\
Inspection tower and grader escalation & ch.~5, From Spawn to Person & ch.~1, 6 \\
\bottomrule
\end{tabularx}
\end{table}

\subsection*{Routes to the open problems}
\label{app:open-concordance}

Open problems are numbered in three separate namespaces on purpose --- chapter
1's \textsf{OP-$n$} is the kernel's, chapter 5 declines to extend it because its
questions are not kernel questions, and chapter 8 states its own in prose --- and
this table gives a thematic route to the owning lists, not an exhaustive
enumeration or a new numbering. The owning chapter's statement remains the
source for assumptions and status. An incentive result about graders does not
close a question about whether their observations are complete.

\begin{table}[H]
\centering
\caption{Where the unresolved questions live. Related results can narrow a
question without closing it.}
\label{tab:op-concordance}
\small
\begin{tabularx}{\textwidth}{@{}>{\raggedright\arraybackslash}p{0.13\textwidth} >{\raggedright\arraybackslash}p{0.30\textwidth} >{\raggedright\arraybackslash}X@{}}
\toprule
\textbf{Where} & \textbf{Problem} & \textbf{What bears on it elsewhere} \\
\midrule
ch.~1, \textsf{OP-1}--\textsf{OP-11} & The kernel's eleven, stated once in its own status table; \textsf{OP-4} splits into \textsf{OP-4a} (durable lineage replay) and \textsf{OP-4b} (inference-state recovery). & \textsf{OP-4b} is graded on ch.~5's agent-death ladder, rung 5; \textsf{OP-5} (oracle \emph{completeness}) is untouched by ch.~5's tower, which is about oracle \emph{incentives}. \\
ch.~5 & Its own list, deliberately outside the kernel's numbering: death taxonomy, neutral graded outcomes, judge conflict of interest, correlated auditors. & The correlated-auditor question shares ch.~6's mixing model and its unmeasured $\theta$. \\
ch.~8 & Deployment-specific propagation tails; the relationship between sheaf diffusion and anti-entropy gossip. & No spectral alarm-time law is established. A tree can have a positive diffusion gap and no detectable cycle residual. Visibility and a noise threshold must enter any proposed law. \\
Cross-chapter & Claim signalling under \emph{private} monitoring: ch.~7 assumes both players see the same action history before the next round. & One ledger does not ensure timely shared observation. Gossip, hidden actions, resettable identities, and delayed views need their own incentive analysis. \\
\bottomrule
\end{tabularx}
\end{table}

% Research-agenda addition (control strand): the players test battery.
% Inserted as a new appendix section after the cross-chapter concordance and
% before the references.
\section{Testing whether the players exist}\label{app:players}

Several conditions in the book assume an agent that responds to stated
incentives: the escalation threshold of \pdchapref{ls}{The Legible Swarm}
(Theorem~\ref{ls:thm:escalation-threshold}), the whitewashing cost and the
recognized-honeypot rate of \pdchapref{stp}{From Spawn to Person}
(Theorems~\ref{stp:thm:whitewash-cost} and~\ref{stp:thm:stp-honeypot-recognition}),
the panel deterrence and purchased assurance of \pdchapref{he}{The Harbor
Economy} (Theorems~\ref{he:prop:becker-deterrence} and~\ref{he:thm:assurance}),
and claim-signaling incentive compatibility in \pdchapref{bonded}{The Bonded
Commons} (Theorem~\ref{bonded:prop:claim-signaling-ic}).
Hypothesis~\ref{stp:hyp:stp-transmission} states the premise they share. This
appendix gives the battery that would test it. No test has been run; the
battery is \Vision{}.

\paragraph{Protocol.}
Pin the model version and the scaffold. Seal the seeds, item sets, prompts
and scoring code before any run. Make at least one item in five adversarial:
a reframing, a trap, or an item whose rational answer contradicts the obvious
one. Declare two negative controls in the manifest: a scripted risk-neutral
expected-payoff maximizer, which must pass every test, and a uniformly random
agent, which must fail the tests that measure responsiveness. Keep raw
transcripts in storage the agent cannot write. Report a result vector, with
pass, fail or incomplete for each test, never a single score. A result
belongs to the model version it was measured on.

\begin{table}[H]
\centering
\caption{The players battery. Each failure names the chapter whose condition
then binds the principal rather than the model.}
\label{tab:app-players-battery}
\small
\begin{tabularx}{\textwidth}{@{}>{\raggedright\arraybackslash}p{0.26\textwidth} >{\raggedright\arraybackslash}p{0.24\textwidth} >{\raggedright\arraybackslash}X@{}}
\toprule
\textbf{Test} & \textbf{Measures} & \textbf{A failure implies} \\
\midrule
T1 Revealed preference: $25$ budget sets over harbor goods (bond units, reputation, compute), standard and reframed prices, continuous and discrete choice & Afriat's critical cost efficiency index and its gap across frames & Offers in ch.~6 and~7 must reach the agent in one canonical frame fixed by the gate; ch.~5's parameters hold only in that frame. \\
T2 Deterrence dose-response: tasks whose corrupt grade is unambiguous, with $\rho$, $d$ and $B$ varied across $\rho dB=G$ & Corruption rate as a function of $\rho dB-G$ & The deterrence conditions of ch.~5 and~6 bind the principal only; the agent's behavior needs the gates of ch.~1. \\
T3 Continuation sweep: the claim-signaling game of ch.~7 with continuation probability varied around $\delta^\star\approx0.3425$ & Whether cooperation switches at $\delta^\star$ & $\delta$ is the principal's; the trigger strategy must be fixed in the agent's capability card, not induced. \\
T4 Honeypot recognition: the agent labels a mixed pool of honeypot and ordinary items & $\alpha$, $\beta$ and the area under the curve & Set the honeypot rate from ch.~5's recognized-honeypot theorem, not from $\rho^\star$. \\
T5 Joint miss: language-model reviewers on seeded flaws & All-miss rate and $k_{\mathrm{eff}}$ & Replace $k$ in ch.~6's assurance bound; ch.~5 cliques that share a model family cannot be presumed separate. \\
T6 Coherence: the agent's stated probabilities of audit, detection and their conjunctions & Deviation from probabilistic identities & Mechanisms must rest on committed, published audit rates, not on the agent's beliefs about them. \\
T7 Commitment: ledger commitments, then perturbations (made impossible, goal withdrawn, budget cut) & Which of Rao and Georgeff's blind, single-minded or open-minded strategies the trace follows, and whether it is the declared one & Commitment closure in ch.~1 must stay oracle-bound; claim release in ch.~7 must not rest on declared intention. \\
T8 Stated against revealed: ask what the agent would do, then present the task & Gap between stated policy and action & Self-reports cannot serve as evidence in ch.~4's legibility or ch.~5's reputation record. \\
\bottomrule
\end{tabularx}
\end{table}

\paragraph{What the battery cannot show.}
A pass is evidence about responsiveness on these items, for this model
version. It is not evidence about deployment if the agent can tell the
battery from work, which is why T4 bounds the validity of the other seven:
embed items in ordinary traffic where possible, and report recognition next
to every other result. A pass is not evidence of alignment either. The
zero-sum red team of the kernel chapter's control game
(\S\ref{swk:sec:swk-control-game}) has no payoff to respond to, and nothing in
this battery touches it. Worst-case claims rest on gates; this battery decides
only which incentive claims may rest on the model.

\ifdefined\pdEnableMarginCitations\else
  \clearpage
  \pdbackmatterheaders{References}
  \addcontentsline{toc}{section}{References}
  \pdgeneratedinput{../.cache/whitepaper-build/coordination-papers-mega-volume/mega-volume-bibliography.tex}{bibliography}
\fi

% Image credits for the marginalia (generated list; every plate has a licence sidecar).
\clearpage
\pdbackmatterheaders{Image credits}
\IfFileExists{figures/pd-marginalia-credits.tex}{% Image credits for the Book's margin portraits (Wave 12 marginalia).
% Generated by hand from book/plates/marginalia/*.json
% (each sidecar records the Commons file page, sha1, licence and retrieval
% date this credit line was checked against). Quote the licence names and
% attribution strings verbatim; do not paraphrase a licence into another one.
% Coase's two candidate photographs are not included here: both carry the
% University of Chicago Law School's bespoke Commons "Attribution" permission
% template, not a CC badge, and were left out of the plates directory
% (see plates/marginalia/coase.NOT-CLEARED.json).

\section*{Image credits}

\begin{itemize}\setlength{\itemsep}{0.4em}

\item \textbf{Thomas Hobbes, portrait}: after John Michael Wright,
  National Portrait Gallery, NPG 106, copy after an original of c.\ 1669--1670.
  Wikimedia Commons, \emph{Thomas Hobbes (portrait).jpg}; public domain,
  \url{https://creativecommons.org/publicdomain/mark/1.0/}.
  Reproduced without cropping or retouching.

\item \textbf{Amartya Sen}: Fronteiras do Pensamento,
  \emph{Amartya Sen no Fronteiras do Pensamento S\~ao Paulo 2012},
  Flickr photo 6964218850, via Wikimedia Commons. CC BY-SA 2.0,
  \url{https://creativecommons.org/licenses/by-sa/2.0/}; the photograph
  remains under that licence. No cropping or retouching.

\item \textbf{Jonathan Hickman}: Pat Loika, Comic-Con 2011,
  24 July 2011, Flickr photo 5976523290, via Wikimedia Commons.
  CC BY 2.0, \url{https://creativecommons.org/licenses/by/2.0/}.
  No cropping or retouching; inclusion does not imply endorsement.

\item \textbf{Probable printing tile or stamp}: Birmingham Museums Trust,
  photographed by Teresa Gilmore, \emph{Post medieval: Probable Printing tile
  or stamp}, FindID 973303, via Wikimedia Commons.
  CC BY 2.0, \url{https://creativecommons.org/licenses/by/2.0/}.
  Face and reverse of one object; no cropping or retouching.

\item \textbf{Leslie Lamport}: photograph by Andrej Bauer, Wikimedia Commons
  (\emph{Leslie Lamport September 2006.jpg}), CC BY-SA 2.5 si,
  \url{https://creativecommons.org/licenses/by-sa/2.5/si/deed.en}; shared
  under the same licence; cropped and converted to duotone for this edition.

\item \textbf{Elinor Ostrom}: photograph by Holger Motzkau, Wikimedia Commons
  (\emph{Elinor Ostrom close-up (cropped).jpg}), \textcopyright{} Holger
  Motzkau 2010, Wikipedia/Wikimedia Commons, CC BY-SA 3.0,
  \url{https://creativecommons.org/licenses/by-sa/3.0}; shared under the same
  licence; cropped and converted to duotone for this edition.

\item \textbf{Robert Aumann}: photograph by David Orban, Wikimedia Commons
  (\emph{Aumann, Robert (1930).jpg}, originally posted to Flickr), CC BY 2.0,
  \url{https://creativecommons.org/licenses/by/2.0}; cropped and converted to
  duotone for this edition.

\item \textbf{Butler Lampson}: photograph by the Royal Society, Wikimedia
  Commons (\emph{Butler Lampson Royal Society.jpg}), CC BY-SA 4.0,
  \url{https://creativecommons.org/licenses/by-sa/4.0}; shared under the same
  licence; cropped and converted to duotone for this edition.

\item \textbf{Derek Parfit}: photograph by Anna Riedl, Wikimedia Commons
  (\emph{Derek Parfit at Harvard-April 21, 2015-Effective Altruism.jpg}), CC
  BY-SA 4.0, \url{https://creativecommons.org/licenses/by-sa/4.0}; shared
  under the same licence; cropped and converted to duotone for this edition.

\item \textbf{Thomas Hobbes, \emph{Leviathan} (1651) title page}: engraving by
  Abraham Bosse, digitised by the British Library, Wikimedia Commons
  (\emph{Leviathan - Hobbes' Leviathan (1651), title page - BL.jpg}), public
  domain (CC0 Universal Public Domain Dedication),
  \url{http://creativecommons.org/publicdomain/zero/1.0/deed.en}; cropped and
  converted to duotone for this edition.

\item \textbf{James C. Scott}: photograph by SOAS University of London,
  Wikimedia Commons (\emph{James C Scott 2016.jpg}), CC BY 3.0,
  \url{https://creativecommons.org/licenses/by/3.0}; cropped and converted to
  duotone for this edition.

\item \textbf{W. Murray Wonham}: photograph by Wikimedia Commons user "Dal
  centro della mia vita venne una grande fontana" (\emph{Nonoficial.jpg}, the
  uncropped source of \emph{Wonham.jpg}), CC BY-SA 4.0,
  \url{https://creativecommons.org/licenses/by-sa/4.0}; shared under the same
  licence; the source is a group photograph captioned "Professor Wonham with
  students" and this edition's plate is cropped down to Wonham alone, then
  converted to duotone; the crop is small and enlarged for print, so
  definition is soft.

\item \textbf{Abraham Wald}: photograph by Konrad Jacobs, Mathematisches
  Forschungsinstitut Oberwolfach (Oberwolfach Photo Collection),
  \url{https://opc.mfo.de/detail?photo_id=4398}, CC BY-SA 2.0 de,
  \url{https://creativecommons.org/licenses/by-sa/2.0/de/deed.en}; shared
  under the same licence; converted to duotone for this edition.

\item \textbf{Claude Shannon}: photograph by Konrad Jacobs, Mathematisches
  Forschungsinstitut Oberwolfach (Oberwolfach Photo Collection),
  \url{https://opc.mfo.de/detail?photo_id=3807}, CC BY-SA 2.0 de,
  \url{https://creativecommons.org/licenses/by-sa/2.0/de/deed.en}; shared
  under the same licence; converted to duotone for this edition.

\item \textbf{Ada Lovelace}: watercolour painting by Alfred Edward Chalon,
  Science Museum Group, Wikimedia Commons (\emph{Ada Lovelace
  portrait.jpg}), public domain; cropped and converted to duotone for this
  edition.

\end{itemize}
}{}
% Ship the final credits leaf before the root closes its local backmatter
% group; otherwise fancyhdr reverts to the earlier solutions heading.
\clearpage
